\subsubsection{1C. The junta citation failure and what OZ actually says}\label{c.-the-junta-citation-failure-and-what-oz-actually-says}

An earlier claimed full Dobrushin proof using ``O'Donnell--Zhao 2016 Theorem 3.5 with modified log-Sobolev parameter \(\rho\)'' is \textbf{withdrawn}.

Checked source: O'Donnell--Zhao, \emph{Polynomial bounds for decoupling, with applications}, \href{https://arxiv.org/html/1512.01603v1}{arXiv:1512.01603v1}:

\begin{itemize}
\tightlist
\item
  There is no Theorem 3.5 and no statement of a general-measure MLSI junta theorem with \(\rho\).
\item
  \textbf{Corollary 3.5} concerns uniform Boolean inputs, assumes a small maximum influence of order \(\exp(-Ct^2d^2)\), and uses Gaussian invariance.
\item
  \textbf{Corollary 3.6} is a uniform-measure Fourier-tail theorem: if \(f\) is not an \((\varepsilon,2^{O(k^2/\varepsilon)})\)-junta then a high-level tail has size \(>\exp(-O(k^2)/\varepsilon)\). For bounded degree \(d\), this gives the usual uniform-product junta approximation.
\item
  Ingredients include Walsh orthogonality, commuting independent noise, Gaussian one-block decoupling, and the MOO invariance principle. Glauber MLSI alone does not supply them.
\end{itemize}

\textbf{Do not revive that citation.} Junta approximation over correlated \(\mu\) from a Glauber semigroup MLSI constant remains unproved. ArXiv/DOI links and the checked text are recorded in \texttt{notes/dfko-repair-outcomes.md} and \texttt{.scratch/chk/oz.txt}.
